\chapter{Rigor metodológico}
\label{cap:rigor}

% BLOQUE «Planteamiento de la comparativa» (2/3). Este capitulo es lo que hace que las
% medidas del bloque sean fiables: sin el, la comparativa se apoyaria en etiquetas
% contaminadas.

% GUION. CAPITULO DIFERENCIAL: es un resultado propio, no relleno.
% Fuente: TFM-Quantum/doc/changes/2026-08_regeneracion_v2.md (registro completo)
%         doc/memoria/materiales.md seccion 3
% Tono: honesto y sin dramatismo. Un trabajo que encuentra y corrige sus propios errores
% vale mas que uno donde todo salio a la primera.

\section{Por qué este capítulo existe}
\label{sec:porque_rigor}
Durante el análisis exploratorio se descubrió que \textbf{las etiquetas del conjunto de datos
estaban contaminadas}. Se midió el alcance, se corrigió y se regeneró. Documentarlo no es un
gesto de transparencia: la razón por la que el error \textbf{pasó inadvertido} es un resultado
metodológico aprovechable.

\section{El error: degradación silenciosa del simulador}
\label{sec:bug}
El estimador ruidoso no fijaba el método de simulación, así que Qiskit Aer usaba
\texttt{automatic} y elegía \textbf{según la memoria libre de la máquina}: matriz de densidad
—exacta— mientras cupiera, y \textbf{trayectorias de Kraus} —Montecarlo— en cuanto dejaba de
caber. Sin ningún aviso.

En una máquina de 11\,GB el cambio ocurría en $n = 10$.

\subsection{Alcance medido}
% Tabla: ver materiales.md seccion 3
Ruido introducido $\sigma \approx 1{,}5\times10^{-3}$, frente a un $|\Delta|$ típico de
0,007--0,019. Afectaba al \textbf{33,9\,\%} del conjunto y al \textbf{56,3\,\%} del test.

\subsection{Por qué importaba, y por qué no tanto como parece}
Para el $R^2$ no era grave: el techo seguía en 0,98--0,997. \textbf{Para el MAE sí}: imponía
un suelo de 0,0012, el \textbf{40\,\%} del listón en el observable más pequeño. Ni un modelo
perfecto habría bajado de ahí.

\section{La lección: qué demuestra realmente una comprobación}
\label{sec:leccion}
% ESTA es la seccion que da valor al capitulo. Redactarla con cuidado.
La documentación afirmaba que las etiquetas eran analíticas, \textit{«verificado ejecutando el
mismo circuito 5 veces con resultado idéntico bit a bit»}.

\textbf{Esa comprobación no demostraba lo que decía demostrar.} Repetir con la \textbf{misma
semilla} prueba \textbf{reproducibilidad}, no ausencia de ruido: un simulador estocástico con
semilla fija devuelve siempre el mismo número, y ese número puede estar igual de equivocado
las cinco veces.

\subsection{Y el mismo error se repitió al investigarlo}
La primera comprobación con semillas \textbf{distintas} se hizo en $n = 8$, donde el método
todavía era exacto, y dio «sin dispersión». Solo al repetirla en $n = 10$ y $n = 12$ apareció
el problema.

\subsection{La regla que queda}
Para afirmar que un cálculo es exacto hay que variar \textbf{la fuente de aleatoriedad}, y
hacerlo \textbf{en todo el rango de parámetros}, no en un punto cómodo.

\section{La corrección}
\label{sec:correccion}
\begin{itemize}
    \item Método \textbf{forzado} a matriz de densidad: si no cabe en memoria,
          \textbf{falla ruidosamente} en vez de degradar la exactitud en silencio.
    \item Cada muestra registra el método usado. \textbf{Es lo que habría delatado el error en
          dos segundos.}
    \item Un manifiesto con versiones, semilla y configuración.
\end{itemize}

\textbf{Consecuencia operativa:} con matriz de densidad forzada, 15 qubits piden 17,2\,GB por
proceso. Las muestras de ese tamaño \textbf{no se pueden generar en la máquina de desarrollo}
—no es lentitud, es imposibilidad— y hubo que generarlas en la nube.

\section{Validación de la corrección}
\label{sec:validacion}
La firma esperada, y que efectivamente aparece: el valor \textbf{exacto no cambia en ninguna
muestra} —nunca fue estocástico—, las diferencias aparecen \textbf{solo} en $n \geq 10$, y su
magnitud coincide con el $\sigma$ medido.

\section{Un segundo límite: el transpilador no es determinista}
\label{sec:transpilador}
% Fuente: doc/changes/2026-08_regeneracion_v2.md seccion 3.5
Con la semilla del transpilador fijada, \texttt{transpile()} \textbf{no produce siempre el
mismo circuito}: medido sobre 13.217 muestras, afecta al \textbf{12,65\,\%}, con una
divergencia en el valor ruidoso de hasta $1{,}45\times10^{-2}$.

% ⚠️ NO citar las cifras de la primera estimacion (4 % y 1,6e-4): salian de una prueba
% de 24 muestras y se quedaban cortas por un factor de 3 y de 90 respectivamente.

\textbf{Es benigno, y la explicación es física:} el valor exacto sale \textbf{idéntico bit a
bit} —los dos circuitos implementan la misma unitaria— y solo el ruidoso difiere, porque
\textbf{el ruido no conmuta con la unitaria}.

\textbf{Lo que hay que matizar:} el conjunto es reproducible desde las semillas para el
\textbf{plan}, pero \textbf{no al bit} para el circuito transpilado.

\section{Dos trampas de medición detectadas}
\label{sec:trampas}
% Fuente: doc/memoria/materiales.md seccion 5
\subsection{Confundir el eje temporal con el eje de tipo}
Sobre el conjunto completo, la correlación del error con el día es $-0{,}26$/$-0{,}39$ y
parece deriva fuerte. \textbf{Dentro del grupo de entrenamiento cae a $-0{,}03$/$-0{,}08$}:
casi todo era el eje de tipo disfrazado, porque los días tardíos pertenecen solo al test.

\subsection{Soporte escaso disfrazado de tendencia}
Una curva por número de qubits parecía caer a partir de 10 qubits. El tramo se sostenía sobre
73, 3 y 2 muestras. \textbf{Publicar siempre el tamaño de la muestra.}
